%%%%%%%%%%%%%%%%%%%%%%%%%%%%%%%%%%%%%%%%%
% @Author: Adam Verner
% Note:
% The \lipsum[#] commands throughout this template generate dummy text
% to fill the template out. These commands should all be removed when 
% writing essay content.
%
%%%%%%%%%%%%%%%%%%%%%%%%%%%%%%%%%%%%%%%%%

%----------------------------------------------------------------------------------------
%	PACKAGES AND OTHER DOCUMENT CONFIGURATIONS
%----------------------------------------------------------------------------------------

\documentclass[12pt,a4paper]{article} % Default font size is 12pt, it can be changed here

% right should be 2, but that behaves weridly
% copied from LibreOffice
\usepackage[left=3cm, right=3cm, top=-1cm, bottom=2.5cm,includehead,includefoot]{geometry} 



% header & footer START
\usepackage{fancyhdr}
\pagestyle{fancy}
\lhead{Střední průmyslová škola na Proseku}
\rhead{26-41-M/01 Elektrotechnika}
\renewcommand{\headrulewidth}{0.4pt}
% header & footer STOP


\usepackage{helvet}
\renewcommand{\familydefault}{\sfdefault}
\renewcommand{\baselinestretch}{1.5} % line spacing

\usepackage[nottoc,numbib]{tocbibind}  % include bibliography


\usepackage{hyperref}

\usepackage{geometry} % Required to change the page size to A4
\geometry{a4paper} % Set the page size to be A4 as opposed to the default US Letter

\usepackage{graphicx} % Required for including pictures

\usepackage{float} % Allows putting an [H] in \begin{figure} to specify the exact location of the figure
\usepackage{wrapfig}
\usepackage{adjustbox}
\usepackage{xcolor}

\usepackage{lipsum} % Used for inserting dummy 'Lorem ipsum' text into the template

\usepackage{mathptmx}
\usepackage{anyfontsize}
\usepackage{t1enc}


\setlength\parindent{35pt}  % odrazeni odstavce

\graphicspath{{Pictures/}} % Specifies the directory where pictures are stored

% use packages to support czech
\usepackage[utf8]{inputenc}
\usepackage[czech]{babel}

\begin{document}

%----------------------------------------------------------------------------------------
%	TITLE PAGE
%----------------------------------------------------------------------------------------

\begin{titlepage}

\begin{minipage}{0.15\textwidth}
	\includegraphics[height=\textwidth]{Praha_logo.eps}
\end{minipage}
\begin{minipage}{0.3\textwidth}
	\includegraphics[height=\textwidth]{SPS-Prosek_logo.jpg}
\end{minipage}
\begin{adjustbox}{minipage=0.55\textwidth,left}
	\textbf{\textcolor{darkgray}{
	Střední průmyslová škola na Proseku\\
	90 00 Praha 9, Novoborská 2}}
\end{adjustbox}

\center

\textbf{\large OBOR}\\[0.5cm]
\textbf{\large	26-41-M/01 ELEKTROTECHNIKA\\
				(ŠVP: Průmyslová automatizace a inteligentní budovy)}\\[4.5cm]
\textbf{\LARGE MATURITNÍ PROJEKT}\\[1.5cm]
\textsc{TÉMA PRÁCE}\\[1.5cm]
\textbf{\color{red} \huge \bfseries  TÉMA PROJEKTU (ČERVENÝ TEXT NAHRAĎTĚ)}

\vspace*{\fill}

\begin{minipage}{0.49\textwidth}
\begin{flushleft} \large
\the\year % insert current year
\end{flushleft}
\end{minipage}
\begin{minipage}{0.49\textwidth}
\begin{flushright}
\textcolor{red}{\large AUTOR (PŘÍJMENÍ Jméno)}
\end{flushright}
\end{minipage}


\end{titlepage}


\newgeometry{left=3cm, right=3cm, top=1.25cm, bottom=2.5cm,includehead,includefoot}


%----------------------------------------------------------------------------------------
%	LIST SE ZADÁNÍM
%----------------------------------------------------------------------------------------
\textbf{\color{red} \LARGE MÍSTO TOHOTO LISTU VLOŽTE LIST SE ZADÁNÍM}
\newpage

%----------------------------------------------------------------------------------------
%	PROHLÁŠENÍ
%----------------------------------------------------------------------------------------
\section*{PROHLÁŠENÍ}

{Prohlašuji, že jsem svou maturitní práci vypracoval samostatně a použil jsem pouze podklady (literaturu, projekty, SW, atd.) uvedené v seznamu.
}\\

{Nemám závažný důvod proti užití tohoto školního díla ve smyslu § 60 zákona č. 121/2000 Sb., o právu autorském, o právech souvisejících s právem autorským a o změně některých zákonů (autorský zákon).}\\[0.5cm]

\begin{minipage}{0.49\textwidth}
\begin{flushleft}
{V Praze dne \today}\\% insert current year
\bigskip
\bigskip
\end{flushleft}
\end{minipage}
\begin{minipage}{0.49\textwidth}
\begin{flushright}
\centering
{
..........................................\\
podpis
}
\end{flushright}
\end{minipage}

\newpage
%----------------------------------------------------------------------------------------
%	PODĚKOVÁNÍ
%----------------------------------------------------------------------------------------
\textsc{\textbf{\fontsize{20pt}{0}\selectfont poděkování}}\\

{\color{red} Text poděkování …. Běžně se děkuje rodinným příslušníkům za jejich morální podporu, vedoucí práce za její vedení, dalším konkrétním učitelům a konzultantům za věcné připomínky atd.}

\newpage
%----------------------------------------------------------------------------------------
%	ANOTACE
%----------------------------------------------------------------------------------------
\section*{ANOTACE}

\begin{table}[htb]

	\begin{tabular}{|p{0.3\linewidth}|p{0.65\linewidth}|} \hline
	Jméno autora          & \\ \hline
	Název maturitní práce & {\color{red} Musí se shodovat s oficiálním názvem práce v zadání}										\\ \hline
	Rozsah práce          & {\color{red} Počet stran práce ve formátu (např. 165 s. + 50 s. příloh)}									\\ \hline
	Školní rok vyhotovení & 2019 \\ \hline
	Škola                 & Střední průmyslová škola na Proseku \\ \hline
	Vedoucí práce         & \\ \hline
	Využití práce         & {\color{red} Zde popište použití práce, např. výukové účely v oboru Telekomunikace v předmětu tom a tom}					\\ \hline
	Abstrakt              & {\color{red} Např. Práce je zaměřena na návrh webové aplikace pro evidenci ...} \\ \hline
	Klíčová slova         & {\color{red} Zde uveďte klíčová slova práce, tj. slova, která definují její obsah, např. webové stránky, skriptovací jazyky, syntaxe,...} \\ \hline
	\end{tabular}
\end{table}
\newpage
%----------------------------------------------------------------------------------------
%	ANNOTATIONS
%----------------------------------------------------------------------------------------
\section*{ANNOTATION}

\begin{table}[htb]
	\begin{tabular}{|p{0.3\linewidth}|p{0.65\linewidth}|} \hline
	Autor                    &							\\ \hline
	Title of graduation work &							\\ \hline
	Extend                   & {\color{red} e.g.: 165 p. + 50 p. of appendices}	\\ \hline
	Academic year            & 2019							\\ \hline
	School                   & Střední průmyslová škola na Proseku			\\ \hline
	Supervisor               &							\\ \hline
	Application              &							\\ \hline
	Abstrakt                 &							\\ \hline
	Key words                &                          \\ \hline 
	\end{tabular}
\end{table}


\newpage
%----------------------------------------------------------------------------------------
%	TABLE OF CONTENTS
%----------------------------------------------------------------------------------------

\tableofcontents % Include a table of contents

\newpage

%----------------------------------------------------------------------------------------
%	INTRODUCTIONlipsum
%----------------------------------------------------------------------------------------

\section{Úvod} % Major section

{\color{red} Text}

%------------------------------------------------

\subsection{Část 1} % Sub-section

{\color{red} Název kapitol bez teček a každá na nové stránce.}\\\\
\lipsum[1] % Dummy text
Example citation \cite{1}.

%------------------------------------------------

\subsection{Část 2} % Sub-section

\lipsum[2] % Dummy text

%------------------------------------------------

\subsubsection{Podčást 1} % Sub-sub-section

\lipsum[3] % Dummy text

\begin{figure}[H] % Example image
\center{\includegraphics[width=0.5\linewidth]{placeholder}}
\caption{ukázkový obrázek.}
\label{fig:speciation}
\end{figure}

%------------------------------------------------

\subsubsection{Podčást 2} % Sub-sub-section

\lipsum[4] % Dummy text

%----------------------------------------------------------------------------------------
%	MAJOR SECTION 1
%----------------------------------------------------------------------------------------

\section{Zápletka} % Major section

\lipsum[5] % Dummy text

%------------------------------------------------

\subsection{Část 1} % Sub-section

\subsubsection{Podčást 1} % Sub-sub-section

\lipsum[6] % Dummy text

%------------------------------------------------

\subsubsection{Podčást 2} % Sub-sub-section


%------------------------------------------------

\subsubsection{Podčást  3} % Sub-sub-section

\begin{description} % Numbered list example

\item[První] \hfill \\
\lipsum[9] % Dummy text

\item[Druhý] \hfill \\
\lipsum[10] % Dummy text

\item[Třetí] \hfill \\
\lipsum[11] % Dummy text

\end{description} 

%----------------------------------------------------------------------------------------
%	MAJOR SECTION X - TEMPLATE - UNCOMMENT AND FILL IN
%----------------------------------------------------------------------------------------

%\section{Content Section}

%\subsection{Subsection 1} % Sub-section

% Content

%------------------------------------------------

%\subsection{Subsection 2} % Sub-section

% Content

%----------------------------------------------------------------------------------------
%	CONCLUSION
%----------------------------------------------------------------------------------------

\section{Závěr} % Major section

\lipsum[12]

%----------------------------------------------------------------------------------------
%	BIBLIOGRAPHY
%----------------------------------------------------------------------------------------

\begin{thebibliography}{99} % Bibliography - this is intentionally simple in this template

\bibitem[1]{1} Seznam citací dle ISO 690 (např. dle \url{www.boldis.cz/citace/citace2.pdf})
\bibitem[2]{2} Nebo dle generátoru citací na \url{www.citace.com}
 
\end{thebibliography}

%----------------------------------------------------------------------------------------

\end{document}

